% $Header: /home/vedranm/bitbucket/beamer/solutions/conference-talks/conference-ornate-20min.en.tex,v 90e850259b8b 2007/01/28 20:48:30 tantau $

\documentclass{beamer} %[compress]
% alternatively add handout or draft option

%\mode<handout>%{\setbeamercolor{background canvas}{bg=black!5}}
%\usepackage{pgfpages}
%\pgfpagesuselayout{4 on 1}[a4paper,border shrink=5mm,landscape]
% For handout printing, 4 slides on 1 page

\mode<presentation>
{
  \usetheme{Pittsburgh} % plain template Pittsburgh
  \useoutertheme{infolines} % bottom (author, title, place)
  \useoutertheme[subsection=false]{miniframes} % top (sections, frames)
  \usecolortheme{beaver} %blue color whale
	\setbeamertemplate{frametitle}[default][right] % frame titles on the right
  \setbeamercovered{transparent}
  % or whatever (possibly just delete it)
}

\setbeamersize{text margin left=1cm, text margin right=1cm}
%\setbeamertemplate{navigation symbols}{}
\usepackage[english]{babel}
\usepackage{booktabs}
\usepackage{comment}
% or whatever

\usepackage[cp1250]{inputenc}
% or whatever
\usepackage{subfigure} %subfigures

\usepackage{times}
\usepackage{tikz}
%\usepackage{breqn}
\usepackage{amsmath}
\usepackage[T1]{fontenc}
\def\sym#1{\ifmmode^{#1}\else\(^{#1}\)\fi} %shortcut for Stata tables
% Or whatever. Note that the encoding and the font should match. If T1
% does not look nice, try deleting the line with the fontenc.

%\logo{\includegraphics[height=0.5cm]{logo.jpg}} % logo on every page



\title[Lecture 9: Heterogeneity] % (optional, use only with long paper titles)
{Lecture 9: Heterogeneity}

\author[T.\ Havranek] % (optional, use only with lots of authors)
{ Tomas Havranek
}
% - Give the names in the same order as the appear in the paper.
% - Use the \inst{?} command only if the authors have different
%   affiliation.

\institute[Prague] % (optional, but mostly needed)
{
  	Charles University, CEPR, METRICS

  %\inst%
  %Economic Research Department\\
  %Czech National Bank
      %\inst{1}%
  %Research Department\\
  %Czech National Bank
  %\and
  %\inst{2}%
  %Institute of Economic Studies\\
  %Charles University, Prague
  
%\pgfdeclareimage[height=1cm]{logo}{logo.jpg} % logo only on title page
%\pgfuseimage{logo}
}% - Use the \inst command only if there are several affiliations.
% - Keep it simple, no one is interested in your street address.

\date[Meta in Econ and Finance] % (optional, should be abbreviation of conference name)
{Research Synthesis in Economics and Finance, \\ University of Canterbury\\[2ex]
  March 17, 2025
}%\\[4mm]% - Either use conference name or its abbreviation.
% - Not really informative to the audience, more for people (including
%   yourself) who are reading the slides online


%\subject{Intertemporal Substitution}

\begin{comment}
\AtBeginSection[]
{
  \begin{frame}<beamer>{Outline}
    \tableofcontents[currentsection]%,currentsubsection]
  \end{frame}
}
\end{comment}

\begin{document}

\begin{frame}
  \titlepage
\end{frame}

\begin{comment}

\begin{frame}{Outline}
  \tableofcontents
  % You might wish to add the option [pausesections]
\end{frame}

\end{comment}

%\section{Summary}
%\subsection{}

%\section{EIS}
%\subsection{}


\begin{frame}{Identification of publication bias}

\bigskip

PEESE:
%\smallskip
		   \begin{equation*}
    \hat{E_i} = \tikz[baseline=(E.base)]{\node[draw,circle,inner sep=1pt] (E) {$E_0$};} + \beta SE(\hat{E})^2_i + u_i,
    \end{equation*}
    \begin{tikzpicture}[overlay, remember picture]
        % Arrow from E0 to SE
        \draw[->,thick,bend right=50] (7.3,1.2) to node[above] {\footnotesize{Methods or p-hacking}} (3.32,1.2);
        % Arrow from SE to u
        \draw[->,thick,bend left=50] (7.3,0.5) to node[below] {\footnotesize{Methods or p-hacking}} (5.9,0.5);
    \end{tikzpicture}
		
%ignoring PET and weights for simplicity.

\smallskip
\smallskip
\smallskip
\smallskip

\bigskip


corr($SE,u) \neq 0$ $\Rightarrow$ $\hat{\beta}$ and $\hat{E_0}$ biased. 


\bigskip
\smallskip
\smallskip
\smallskip

%\alert{Natural solution}: ue N as an insturment for $SE(\hat{E})^2_i$ $\rightarrow$ MAIVE. %corr.\ with $SE(\hat{E})^2_i$, not with methods or p-hacking

   \fbox{\parbox{0.99\linewidth}{\alert{Natural solution:} N instrumenting $SE(\hat{E})_i^2 \rightarrow$ MAIVE.}}
	
	\bigskip
	
	Or add covariates!

\end{frame}




\begin{frame}{Classical heterogeneity measure}
 
    
    \textbf{Common measure:} estimate the relative heterogeneity of underlying effects \textit{(Higgins and Thompson, 2002)}
    
    \[
    I^2 \equiv \frac{\text{underlying variance of effects}}{\text{observed variance of estimates}} = \frac{\tau^2}{\tau^2 + \bar{v}}
    \]
    
    \begin{itemize}
        \item $I^2 = 0.75 \sim 1.0$ \dots \textit{``considerable heterogeneity'' in medicine (Cochrane Review Guideline)}
        \item $I^2 \simeq 0.9$ \dots common in economics meta-analyses (e.g., \textit{Vivalt, 2021})
    \end{itemize}
    
    \bigskip
    
    \textbf{Unresolved concern:}
    \begin{itemize}
        \item no universally agreed criteria for ``too much'' heterogeneity
        \item undesirable pattern: with more precise studies, higher $I^2$
    \end{itemize}
\end{frame}




\begin{frame}{Guiding example: beauty and success}
 \begin{figure}
	 \centering
		 \includegraphics[width=0.9\textwidth]{1.png}
 \end{figure}
\end{frame}

\begin{frame}{Heterogeneity both within and across}
 \begin{figure}
	 \centering
		 \includegraphics[width=1.00\textwidth]{2.png}
 \end{figure}
\end{frame}

\begin{frame}{Wide dispersion}
 \begin{figure}
	 \centering
		 \includegraphics[width=0.9\textwidth]{3.png}
 \end{figure}
\end{frame}

\begin{frame}{Measurement of beauty}
 \begin{figure}
	 \centering
		 \includegraphics[width=1.00\textwidth]{4.png}
 \end{figure}
\end{frame}

\begin{frame}{Measurement of success}
 \begin{figure}
	 \centering
		 \includegraphics[width=1.00\textwidth]{5.png}
 \end{figure}
\end{frame}

\begin{frame}{Data characteristics}
 \begin{figure}
	 \centering
		 \includegraphics[width=1.00\textwidth]{6.png}
 \end{figure}
\end{frame}

\begin{frame}{Estimation technique}
 \begin{figure}
	 \centering
		 \includegraphics[width=1.00\textwidth]{7.png}
 \end{figure}
\end{frame}

\begin{frame}{Publication characteristics}
 \begin{figure}
	 \centering
		 \includegraphics[width=1.00\textwidth]{8.png}
 \end{figure}
\end{frame}

\begin{frame}{Beauty measurement doesn't matter}
 \begin{figure}
	 \centering
		 \includegraphics[width=0.8\textwidth]{9.png}
 \end{figure}
\end{frame}

\begin{frame}{Success measurement doesn't matter}
 \begin{figure}
	 \centering
		 \includegraphics[width=0.8\textwidth]{10.png}
 \end{figure}
\end{frame}

\begin{frame}{Method doesn't matter}
 \begin{figure}
	 \centering
		 \includegraphics[width=0.8\textwidth]{11.png}
 \end{figure}
\end{frame}


\begin{frame}{Ability control matters}
 \begin{figure}
	 \centering
		 \includegraphics[width=0.8\textwidth]{12.png}
 \end{figure}
\end{frame}


\begin{frame}{Occupation matters}
 \begin{figure}
	 \centering
		 \includegraphics[width=0.8\textwidth]{13.png}
 \end{figure}
\end{frame}

\begin{frame}{Sex workers stand apart}
 \begin{figure}
	 \centering
		 \includegraphics[width=0.9\textwidth]{14.png}
 \end{figure}
\end{frame}


\begin{frame}{Option 1: subsamples}
 \begin{figure}
	 \centering
		 \includegraphics[width=1.00\textwidth]{15.png}
 \end{figure}
\end{frame}

\begin{frame}{Option 1: subsamples}
 \begin{figure}
	 \centering
		 \includegraphics[width=1.00\textwidth]{16.png}
 \end{figure}
\end{frame}

\begin{frame}{Option 1: subsamples}
 \begin{figure}
	 \centering
		 \includegraphics[width=1.00\textwidth]{17.png}
 \end{figure}
\end{frame}


\begin{frame}{Option 1: subsamples}
 \begin{figure}
	 \centering
		 \includegraphics[width=1.00\textwidth]{18.png}
 \end{figure}
\end{frame}


\begin{frame}{Option 2: Bayesian model averaging}
\begin{itemize}
	\item We want to regress estimates of the beauty effect on variables reflecting heterogeneity
	\item But too many variables; model uncertainty
	\item Solution: run regressions with different combinations of the variables
	\item Give more weight to those that fit the data well and are parsimonious
\end{itemize}
\end{frame}


\begin{frame}{Option 2: Bayesian model averaging}
\begin{itemize}
\item Model averaging can be frequentist but computationally difficult
\item Intuition: run different combinations and weight by adjusted $R^2$ or information criteria
\item In the Bayesian setting we can use the MCMC algorithm
\item G-prior: all coefficients are zero. Weight? UIP = the prior has the same weight as one observation
\item Model prior: uniform = all models have the same weight
\item Dilution prior: models with collinearity get less weight
\end{itemize}
\end{frame}



\begin{frame}{Model inclusion in BMA}
 \begin{figure}
	 \centering
		 \includegraphics[width=0.7\textwidth]{19.png}
 \end{figure}
\end{frame}

\begin{frame}{Posterior densities}
 \begin{figure}
	 \centering
		 \includegraphics[width=0.95\textwidth]{posterior.png}
 \end{figure}
\end{frame}


\begin{frame}{Regression results}
 \begin{figure}
	 \centering
		 \includegraphics[width=0.7\textwidth]{20.png}
 \end{figure}
\end{frame}


\begin{frame}{Best practice (implied estimate)}
\begin{itemize}
\item From BMA results we know how estimates depend on study design
\item Some study designs are better: identification, data, generally quality
\item Select best practice based on the literature (quasi-experimental studies, etc.)
\item Compute fitted values and confidence intervals (lincom command in Stata)
\end{itemize}
\end{frame}


\begin{frame}{Best practice (class size effect)}
 \begin{figure}
	 \centering
		 \includegraphics[width=0.8\textwidth]{21.png}
 \end{figure}
\end{frame}


\end{document}


